%
% Copyright © 2018 Peeter Joot.  All Rights Reserved.
% Licenced as described in the file LICENSE under the root directory of this GIT repository.
%
%{

\makeproblem{Linear magnetic density and currents.}{problem:statics:240}{
Given magnetic charge density \( \rho_m = \lambda_m \delta(x) \delta(y) \), and current density \( \BM = v \Be_3 \rho_m = \Bv \rho_m \), show that the field is given by
\begin{equation*}
F(\Bx) = \frac{\lambda_m c}{2 \pi R} I \rhocap \lr{ 1 - \frac{\Bv}{c} },
\end{equation*}
or with \( \BB = \lambda_m \rhocap/(2 \pi R) \),
\begin{equation*}
F = \BB \cross \Bv + c I \BB.
\end{equation*}
} % problem

\makeanswer{problem:statics:240}{
With only magnetic sources, the current density multivector is
\begin{equation}\label{eqn:statics_infiniteLineChargeMagneticLineCharge:20}
\begin{aligned}
J
&= I \lr{ c \rho_\txtm - \BM } \\
&= I c \lr{ 1 - \frac{v}{c} \Be_3 } \lambda_m \delta(x) \delta(y).
\end{aligned}
\end{equation}
As in \cref{eqn:statics_infiniteLineCharge:160}, let the field observation point be \( \Bx = \Bx_\perp + z \Be_3 \), so
\begin{equation}\label{eqn:statics_infiniteLineChargeMagneticLineCharge:40}
\begin{aligned}
F(\Bx)
&= \inv{4\pi} \int_V dV' \frac{\gpgrade{(\Bx - \Bx') I c \lr{ 1 - \frac{v}{c} \Be_3 } \lambda_m \delta(x) \delta(y) }{1,2}}{\Norm{\Bx - \Bx'}^3}  \\
&= \frac{I c \lambda_m}{4\pi} \int_{-\infty}^\infty dz' \frac{\gpgrade{(\Bx_\perp - z' \Be_3) \lr{ 1 - \frac{v}{c} \Be_3 } }{1,2}}{\lr{\Bx_\perp^2 - \lr{z - z'}^2}^{3/2}} \\
\end{aligned}
\end{equation}
The grade selection reduces to three terms
\begin{equation}\label{eqn:statics_infiniteLineChargeMagneticLineCharge:60}
\begin{aligned}
\gpgrade{(\Bx_\perp - z' \Be_3) \lr{ 1 - \frac{v}{c} \Be_3 } }{1,2}
&=
\Bx_\perp - z' \Be_3 - \frac{v}{c} \Bx_\perp \Be_3  \\
&=
\Bx_\perp \lr{ 1 - \frac{v}{c} \Be_3 } - z' \Be_3.
\end{aligned}
\end{equation}
As before the integral with the \(z'\) dependence vanishes, and we are left with
\begin{equation}\label{eqn:statics_infiniteLineChargeMagneticLineCharge:80}
F(\Bx) = \frac{I c \lambda_m \Bx_\perp \lr{ 1 - \frac{v}{c} \Be_3 } }{2\pi \Bx_\perp^2}.
\end{equation}
Substitution of \( \Bx_\perp = R \rhocap \) yields the desired result.

Finally, observe that
\begin{equation}\label{eqn:statics_infiniteLineChargeMagneticLineCharge:100}
\begin{aligned}
I \rhocap \Be_3
&=
I \lr{ \rhocap \wedge \Be_3 } \\
&=
I^2 \lr{ \rhocap \cross \Be_3 },
\end{aligned}
\end{equation}
so
\begin{equation}\label{eqn:statics_infiniteLineChargeMagneticLineCharge:120}
F(\Bx) = I c \BB + \BB \cross \Bv,
\end{equation}
as claimed.

} % answer
%}
